\section{Reasoning y verifier: no fijar las intuiciones demasiado pronto}

La entrevista tuvo lugar poco después de la publicación de DeepSeek R1. La valoración de Lample tiene dos niveles: por un lado, reconoce que un artículo público que presenta experimentos exitosos y fallidos ahorra a toda la comunidad el costo de búsqueda; por otro, advierte que el reasoning sigue en una etapa temprana, y que las intuiciones actuales sobre el tamaño de los modelos, los métodos de RL y el test-time compute podrían volverse obsoletas tan rápido como las intuiciones de la primera época de GPT-3.

\subsection{La unidad de sistema de un reasoning escalable}

El reasoning training necesita un entorno que proporcione tareas y acciones ejecutables, un verifier que juzgue los resultados, una estrategia de muestreo/búsqueda que explore candidatos y, después, usar el reward o el resultado del filtrado para actualizar el modelo. Sin un verifier confiable, una chain-of-thought larga solo aumenta la longitud del texto; no garantiza la corrección.

\lecturefigure{13-reasoning-verifier-loop.png}{La unidad escalable del reasoning es el ciclo cerrado environment--policy--verifier--update.}{Subtítulos locales 27:45--30:45; redibujo conceptual.}

\begin{knowledgebox}{Términos clave: los cuatro niveles de un sistema de reasoning}
\begin{tabularx}{\textwidth}{L{0.18\textwidth}X X}
\toprule
Nivel & Pregunta central & Fallas comunes \\
\midrule
Environment & Si el estado, las acciones y las condiciones de terminación son ejecutables & Tareas que no pueden reiniciarse, herramientas inestables, filtración de estado \\
Verifier & Si puede juzgar el resultado final o los intermedios & Reward hacking, pruebas erróneas, pruebas que solo verifican el formato \\
Search / sampling & Cómo asignar el test-time compute & Candidatos muy correlacionados, costo fuera de control, poda prematura \\
Training update & Cómo mejorar la policy a partir de los resultados & Sobreajuste al verifier, regresión de capacidades, sesgo en los datos \\
\bottomrule
\end{tabularx}
\end{knowledgebox}

\begin{warningbox}{Los experimentos fallidos publicados también son infraestructura de alto valor}
Cuando un artículo explica que PRM, MCTS o cierto diseño de reward no funcionó en su configuración, no se trata de una «noticia negativa», sino de una reducción del espacio de búsqueda de la comunidad bajo condiciones delimitadas. La condición es registrar el entorno, el modelo, los datos y la evaluación; de lo contrario, la falla no puede reproducirse ni trasladarse a otras configuraciones.
\end{warningbox}

\teachervoice{El docente no presentó DeepSeek R1 como un vuelco del roadmap de Mistral, sino que consideró la investigación abierta como un activo reutilizable. Más importante aún, reconoció que las intuiciones de la era del reasoning siguen siendo inestables. Este «no saber» es más adecuado para incluirse en los apuntes que una certeza fabricada a posteriori.}

\subsection{Resumen de la sección}

El centro de un sistema de reasoning no es un texto más largo, sino un entorno ejecutable y un verifier confiable. En la etapa temprana conviene ampliar la cobertura de los experimentos, publicar los resultados negativos y mantener los límites entre objetivos y evidencia, en lugar de idealizar una sola receta.
